\section{Numerical evidence}\label{sec:numerics}

This section compares \cref{thm:main} with exact computations of $V_f(x,y)$. Everything here is evidence and plays no part in the proofs. All tables and plots are regenerated by \nolinkurl{manuscript/reproduce.py} (see \nolinkurl{manuscript/README.md}). That script re-runs the accepted computation from its source code and recompiles its C\texttt{++} part; it also computes the theoretical quantities. Numbers quoted in this section as bounds or ranges (``at most'', ``between'', ``lies in'') are rounded outwards, so that the printed statement holds for the computed values; numbers quoted as approximate values (``about'', ``approximately'', $\approx$) are rounded to the nearest digit.

\subsection{What was computed}\label{sec:num-method}

\emph{Grid.} The three polynomials are $f(k)=k^2+1$, $k^3+2$ and $2k+1$, with weights $f'(k)=2k$, $3k^2$ and $2$. All three satisfy~\eqref{eq:f-hyp}. The values of $x$ are $x=2.5\cdot10^5\cdot2^j$ for $0\le j\le10$, so the largest is $x=2.56\cdot10^8$. For each $x$ and each target ratio $r\in\{0.01,0.02,0.05,0.1,0.2,0.4,0.7,1\}$, $y$ is the largest integer with $f(y)\le rx$. At the two largest $x$ the ratio $r=0.005$ was added as a diagnostic. The quantity reported is the normalised variance
\[
N(x,r)=\frac{V_f(x,y)}{x^{1/2}f(y)^{3/2}},
\]
computed with the actual ratio $f(y)/x$. By \cref{thm:main}(iii), $N(x,r)\to C_f$ when $x\to\infty$ and $f(y)/x\to0$ suitably.

\emph{Method.} The computation (archived numerical snapshot, files \texttt{compute.py} and \texttt{counts.cpp}) is exact in all integer quantities.
\begin{enumerate}[label=(\arabic*)]
\item The squarefree indicator $\mu^2(n)$, $n\le x$, is sieved. All autocorrelations
\[
C(h)=\sum_{1\le n\le x-h}\mu^2(n)\mu^2(n+h),\qquad0\le h\le x,
\]
are obtained from one number-theoretic transform modulo the prime $p_0=3221225473$, with primitive root~$5$. Every $C(h)$ is at most $x<p_0$, so no information is lost.
\item As in the proof of \cref{prop:identity}, $\sum_lQ(x;q,l)^2=Q(x)+2\sum_{m\ge1}C(mq)$ for every modulus $q=f(k)$.
\item The densities are $g(q,l)=\frac6{\pi^2}\tilde g(q,l)$ with an explicit rational $\tilde g(q,l)$ equivalent to \cref{prop:densities}(i),(iii). The coefficients $\sum_lQ(x;q,l)\tilde g(q,l)$ and $\sum_l\tilde g(q,l)^2$ are computed as exact rationals, using the counts of squarefree multiples of $d$ as in the proof of \cref{prop:trivial}.
\item Finally $V(x;q)=\sum_lQ^2-2x\sum_lgQ+x^2\sum_lg^2$, as in~\eqref{eq:id1}, is accumulated with the weights $f'(k)$ in $70$-digit decimal arithmetic.
\end{enumerate}
The snapshot validates the implementation in several ways. It checks the transform against direct lag sums at every $x$. It compares $\tilde g$ with brute-force local densities for all $q\le50$ and with empirical frequencies at $x=10^6$. At $x=257$ it compares every generated modulus with direct residue counts. During the original run, the operator and an integrator also recomputed some rows independently by direct counting (Appendix~\ref{app:provenance}).

\emph{Re-run.} For this write-up the snapshot code was run again, unchanged, after its SHA-256 hashes had been checked against the acceptance record. The C\texttt{++} part was compiled locally with \texttt{clang++ -O3}. All $270$ rows ($90$ per polynomial, including the diagnostic rows) reproduce the archived integer counts exactly. The values of $V_f$ agree to all stored digits (relative differences at most $\RerunMaxRelV$), and all fitted quantities are identical. The full grid needs about $5$\,GB of memory; the running times of the re-run that produced the tables are recorded in \texttt{reproduction.json} (fields \texttt{snapshot\_rerun.seconds} and \texttt{snapshot\_rerun.total\_seconds\_of\_reproduce\_py}).

\subsection{The constants \texorpdfstring{$C_f$}{C\_f}}\label{sec:num-constants}

For the three polynomials the root counts are:
\begin{itemize}
\item $k^2+1$: $\rho(2)=1$, $\rho(4)=0$, and $\rho(p)=\rho(p^2)=1+\bigl(\frac{-1}{p}\bigr)$ for odd $p$;
\item $k^3+2$: $\rho(2)=\rho(3)=1$, $\rho(4)=\rho(9)=0$; for $p\ge5$, $\rho(p^2)=\rho(p)$, with $\rho(p)=1$ if $p\equiv2\pmod 3$, and, if $p\equiv1\pmod3$, $\rho(p)=3$ or $0$ according as $(-2)^{(p-1)/3}\equiv1\pmod p$ or not;
\item $2k+1$: $\rho(2)=\rho(4)=0$, and $\rho(p)=\rho(p^2)=1$ for odd $p$.
\end{itemize}
For primes not dividing the discriminant all roots are simple, and each root modulo $p$ lifts to exactly one root modulo $p^2$, which gives $\rho(p^2)=\rho(p)$. The formulas were checked by brute force modulo $p$ and $p^2$ for all $p<\RhoBound$. None of the three polynomials has a content prime, so $P_0=\emptyset$, and~\eqref{eq:sigma-close} holds for every $p$ with $K=3+2d$. Hence, for $N\ge30$,
\[
\Bigl|\log\prod_{p>N}\sigma_p\Bigr|\le\frac{K}{1-K/N^2}\sum_{p>N}\frac1{p^2}\le\frac{8K}{1-K/N^2}\Bigl(\frac1{N^2}+\frac1{30N}\Bigr).
\]
Here we used that every prime $p>N\ge30$ lies in one of the $8$ reduced residue classes modulo $30$, and that each class contributes at most $N^{-2}+(30N)^{-1}$. \Cref{tab:constants} lists the partial values $\frac23I\prod_{p\le N}\sigma_p$ and the resulting interval for $C_f$ with $N=10^7$. The partial values are rounded to nine decimals; the interval endpoints are rounded outwards, so that each printed interval contains the computed one. The intervals account rigorously for the omitted primes. They do not account for floating-point rounding in the partial products, which are computed in double precision. As a check of that rounding (not a certified bound), the partial products for $N=10^5$ were recomputed in $40$-digit arithmetic; they agree with the double-precision values to a relative difference of at most $\FloatCheck$.

\begin{table}[htbp]
\centering
\footnotesize\setlength{\tabcolsep}{3.5pt}
\begin{tabular}{lrrrrrl}
\toprule
$f$ & $N=10^{3}$ & $N=10^{4}$ & $N=10^{5}$ & $N=10^{6}$ & $N=10^{7}$ & interval for $C_f$ \\
\midrule
$k^2+1$ & 0.138152461 & 0.138088027 & 0.138083068 & 0.138082663 & 0.138082629 & $[0.138082603,\,0.138082655]$ \\
$k^3+2$ & 0.135675596 & 0.135612820 & 0.135607954 & 0.135607555 & 0.135607522 & $[0.135607489,\,0.135607555]$ \\
$2k+1$ & 0.142685343 & 0.142618897 & 0.142613765 & 0.142613346 & 0.142613311 & $[0.142613291,\,0.142613330]$ \\
$t$ (formal) & 0.128543562 & 0.128483702 & 0.128479079 & 0.128478702 & 0.128478670 & $[0.128478652,\,0.128478687]$ \\
Vaughan $c_2$ & 0.128510919 & 0.128481180 & 0.128478873 & 0.128478684 & 0.128478668 & \\
\bottomrule
\end{tabular}

\caption{Partial Euler products $\frac23I\prod_{p\le N}\sigma_p$ and intervals for $C_f$. The row ``$t$ (formal)'' uses $\rho\equiv1$. The last row is Vaughan's constant~\eqref{eq:vaughan-c2} with $C_1$ truncated at $p\le N$; its limit equals that of the row above (\cref{rem:vaughan}).}
\label{tab:constants}
\end{table}

Apart from floating-point rounding, the intervals therefore fix the digits
\[
C_{k^2+1}=\CfQuadCert\ldots,\qquad C_{k^3+2}=\CfCubicCert\ldots,\qquad C_{2k+1}=\CfLinCert\ldots .
\]
The partial products at $N=10^7$ give the approximate values $C_{k^2+1}\approx\CfQuad$, $C_{k^3+2}\approx\CfCubic$ and $C_{2k+1}\approx\CfLin$; the digits beyond those fixed above are only indicative. For $f(t)=t$ the formal value $\frac23I\prod_{p\le 10^7}\sigma_p\approx\CfT$ agrees with Vaughan's $c_2\approx\VaughanCtwo$ (with $C_1$ truncated at the same point), as \cref{rem:vaughan} predicts. The local identity behind that remark, the identities of \cref{lem:local}, and $\zeta(-\frac12)=-\zeta(\frac32)/(4\pi)$ were also verified with computer algebra and $80$-digit arithmetic (\texttt{reproduce.py}, function \texttt{symbolic\_local\_identity}).

\subsection{The data and the predicted constant}\label{sec:num-data}

\Cref{tab:largest} shows $N(x,r)$ at the largest $x$, together with the relative change from $x=1.28\cdot10^8$. The full tables are in Appendix~\ref{app:tables}. Two features are visible.

First, at fixed $r$ the normalised values are stable in $x$. Between the two largest sizes they change by at most $\StabCubic\%$ for $k^3+2$, and by at most $\StabOther\%$ for $k^2+1$ and $2k+1$ (see also \cref{fig:x}). Fits of $\log V_f$ against $\log x$ at fixed target ratio give $\alpha+\beta\approx2$ (\cref{tab:fits}), consistent with the homogeneous main term $x^{1/2}Y^{3/2}$.

Second, $N(x,r)$ depends clearly on $r$. It rises from values between $\NDiagMin$ and $\NDiagMax$ at $r=0.005$ to values between $\NOneMin$ and $\NOneMax$ at $r=1$. A fit of $\log V_f$ against $\log f(y)$ at fixed $x$ therefore gives exponents $\beta$ between $\BetaMin$ and $\BetaMax$ instead of $\frac32$ (\cref{tab:fits}). This does not contradict \cref{thm:main}. Part (iii) concerns $f(y)/x\to0$, and at a fixed ratio $r$ the error term $x^{1/3+\eps}Y^{5/3}$ of part (i) is as large as the main term. The ratio dependence is captured by the exact-shape statement~(ii); see \cref{sec:num-shape}.

\begin{table}[htbp]
\centering
\footnotesize\setlength{\tabcolsep}{3.5pt}
\begin{tabular}{lrrrrrrrrr}
\toprule
$r$ & $0.005$ & $0.01$ & $0.02$ & $0.05$ & $0.1$ & $0.2$ & $0.4$ & $0.7$ & $1$ \\
\midrule
$k^2+1$ & 0.14039 & 0.14086 & 0.14292 & 0.14446 & 0.15098 & 0.15655 & 0.16278 & 0.16993 & 0.17110 \\
\quad change (\%) & $-0.0233$ & $-0.0363$ & $+0.0051$ & $-0.0102$ & $-0.0118$ & $-0.0043$ & $-0.0083$ & $-0.0003$ & $-0.0030$ \\
$k^3+2$ & 0.14076 & 0.13999 & 0.14120 & 0.14391 & 0.14848 & 0.15234 & 0.15550 & 0.16399 & 0.16671 \\
\quad change (\%) & $+0.4079$ & $-0.9050$ & $-0.5298$ & $-0.3279$ & $-0.3250$ & $-0.0824$ & $-0.2807$ & $-0.1486$ & $-0.1688$ \\
$2k+1$ & 0.14484 & 0.14460 & 0.14837 & 0.14934 & 0.15670 & 0.16171 & 0.16245 & 0.18239 & 0.18369 \\
\quad change (\%) & $+0.0003$ & $-0.0000$ & $+0.0002$ & $+0.0005$ & $+0.0002$ & $+0.0001$ & $+0.0001$ & $-0.0000$ & $-0.0000$ \\
\bottomrule
\end{tabular}

\caption{$N(x,r)=V_f(x,y)/(x^{1/2}f(y)^{3/2})$ at $x=2.56\cdot10^8$ for the target ratios $r$ (the column $0.005$ is the diagnostic ratio). The ``change'' rows give the relative change, in percent, from $x=1.28\cdot10^8$ to $x=2.56\cdot10^8$.}
\label{tab:largest}
\end{table}

\emph{Extrapolation to $r\to0$.} The snapshot estimates $\lim_{r\to0}N$ empirically. At the largest $x$ it fits the model $N(r)=C+a\sqrt r+br$ by unweighted least squares to the rows with $r$ at or above a ``stable'' cutoff. The cutoff is the smallest requested ratio such that every row at or above it changes by less than $1\%$ between the two largest $x$; here it is $\StableCutoff$ for all three polynomials. The reported uncertainty is a \emph{heuristic sensitivity envelope}: the largest of twice the formal standard error of the intercept, the spread over neighbouring cutoffs, the shift to the previous $x$, the shifts under the alternative models $C+a\sqrt r$ and $C+a\sqrt r+br+cr^{3/2}$, and the shifts when the fitted window is cut at $r\le0.1$, $0.2$ or $0.4$. It is not a confidence interval and not an error bound, and the model $C+a\sqrt r+br$ itself is not derived from the theorem. \Cref{tab:fits} compares the extrapolated constants with $C_f$. The relative differences are approximately $\RelQuad\%$, $\RelCubic\%$ and $\RelLin\%$ for $k^2+1$, $k^3+2$ and $2k+1$, all within the envelopes. \Cref{tab:sensitivity} shows how strongly the extrapolated value depends on the model. For $k^2+1$ the alternatives range from about $\SensQuadMin$ to $\SensQuadMax$ and bracket $C_f$, so the difference of $\RelQuad\%$ reflects the choice of model rather than the constant. \Cref{fig:ratio} shows the data, the fitted curves and $C_f$.

\begin{table}[htbp]
\centering
\small
\begin{tabular}{lrrrrrr}
\toprule
$f$ & $\alpha+\beta$ (fixed $r$) & $\beta$ (fixed $x$) & $C_{\mathrm{extr}}\pm$ envelope & $C_f$ & $\frac{C_{\mathrm{extr}}-C_f}{C_f}$ & $\frac{C_{\mathrm{extr}}-C_f}{\mathrm{envelope}}$ \\
\midrule
$k^2+1$ & $2.0008\pm0.0006$ & $1.545\pm0.004$ & $0.1345\pm0.0073$ & $0.13808$ & $-2.61\%$ & $-0.50$ \\
$k^3+2$ & $1.9978\pm0.0021$ & $1.538\pm0.004$ & $0.1363\pm0.0024$ & $0.13561$ & $+0.53\%$ & $+0.30$ \\
$2k+1$ & $2.0001\pm0.0001$ & $1.551\pm0.007$ & $0.1417\pm0.0082$ & $0.14261$ & $-0.65\%$ & $-0.11$ \\
\bottomrule
\end{tabular}

\caption{Empirical exponents and extrapolated constants from the snapshot fits, compared with $C_f$ from \cref{tab:constants}. ``$\alpha+\beta$'' is the mean of the eight fixed-ratio slopes of $\log V_f$ against $\log x$ (with the across-ratio standard error); ``$\beta$'' is the fixed-$x$ slope against $\log f(y)$ at the largest $x$ (with its regression standard error). The last column is $(C_{\mathrm{extr}}-C_f)$ divided by the heuristic envelope; a value of absolute size at most $1$ means that $C_f$ lies within the envelope.}
\label{tab:fits}
\end{table}

\begin{table}[htbp]
\centering
\footnotesize
\begin{tabular}{lrrrrrr}
\toprule
$f$ & primary & $C+a\sqrt r$ & $C+a\sqrt r+br+cr^{3/2}$ & $r\le 0.1$ & cutoff $0.005$ & previous $x$ \\
\midrule
$k^2+1$ & 0.13447 & 0.13837 & 0.13714 & 0.14176 & 0.13547 & 0.13451 \\
$k^3+2$ & 0.13632 & 0.13759 & 0.13569 & 0.13856 & 0.13728 & 0.13763 \\
$2k+1$ & 0.14169 & 0.14086 & 0.14145 & 0.14416 & 0.14188 & 0.14169 \\
\bottomrule
\end{tabular}

\caption{Sensitivity of the extrapolated constant: the primary fit, the two alternative models, the window $r\le0.1$, the lower cutoff $0.005$, and the primary fit at the previous size $x=1.28\cdot10^8$.}
\label{tab:sensitivity}
\end{table}

\begin{figure}[htbp]
\centering
\begin{tikzpicture}
\begin{axis}[width=0.82\textwidth,height=0.5\textwidth,xlabel={$\sqrt{r}$, with $r=f(y)/x$},ylabel={$N(x,r)$},xmin=-0.02,xmax=1.02,ymin=0.13,ymax=0.19,
  legend pos=north west,legend cell align=left,legend style={font=\footnotesize},tick label style={font=\footnotesize},
  yticklabel style={/pgf/number format/fixed,/pgf/number format/precision=3}]
\addplot[only marks,mark=*,mark size=1.6pt,blue] table[x=sqrtr,y=N]{generated/fig_ratio_0.dat};
\addlegendentry{$k^2+1$}
\addplot[only marks,mark=square*,mark size=1.6pt,red] table[x=sqrtr,y=N]{generated/fig_ratio_1.dat};
\addlegendentry{$k^3+2$}
\addplot[only marks,mark=triangle*,mark size=2pt,black!60!green] table[x=sqrtr,y=N]{generated/fig_ratio_2.dat};
\addlegendentry{$2k+1$}
\addplot[blue,dashed,forget plot] table[x=sqrtr,y=N]{generated/fig_fit_0.dat};
\addplot[red,dashed,forget plot] table[x=sqrtr,y=N]{generated/fig_fit_1.dat};
\addplot[black!60!green,dashed,forget plot] table[x=sqrtr,y=N]{generated/fig_fit_2.dat};
\addplot[only marks,mark=o,mark size=3pt,blue,thick,forget plot] table[x=sqrtr,y=N]{generated/fig_cf_0.dat};
\addplot[only marks,mark=o,mark size=3pt,red,thick,forget plot] table[x=sqrtr,y=N]{generated/fig_cf_1.dat};
\addplot[only marks,mark=o,mark size=3pt,black!60!green,thick,forget plot] table[x=sqrtr,y=N]{generated/fig_cf_2.dat};
\end{axis}
\end{tikzpicture}
\caption{Normalised variance at $x=2.56\cdot10^8$ against $\sqrt r$ (filled marks), the snapshot's empirical fits $C+a\sqrt r+br$ (dashed), and the proved constants $C_f$ (open circles at $r=0$).}
\label{fig:ratio}
\end{figure}

\begin{figure}[htbp]
\centering
\begin{tikzpicture}
\begin{axis}[width=0.6\textwidth,height=0.42\textwidth,xmode=log,xlabel={$x$},ylabel={$N(x,r)$},ymin=0.135,ymax=0.19,
  legend pos=outer north east,legend cell align=left,legend style={font=\footnotesize},tick label style={font=\footnotesize},
  yticklabel style={/pgf/number format/fixed,/pgf/number format/precision=3}]
\addplot[blue,mark=*,mark size=1.3pt] table[x=x,y=N]{generated/fig_x_0a.dat};\addlegendentry{$k^2+1$, $r=0.01$}
\addplot[blue,mark=o,mark size=1.3pt,dashed] table[x=x,y=N]{generated/fig_x_0b.dat};\addlegendentry{$k^2+1$, $r=1$}
\addplot[red,mark=square*,mark size=1.3pt] table[x=x,y=N]{generated/fig_x_1a.dat};\addlegendentry{$k^3+2$, $r=0.01$}
\addplot[red,mark=square,mark size=1.3pt,dashed] table[x=x,y=N]{generated/fig_x_1b.dat};\addlegendentry{$k^3+2$, $r=1$}
\addplot[black!60!green,mark=triangle*,mark size=1.6pt] table[x=x,y=N]{generated/fig_x_2a.dat};\addlegendentry{$2k+1$, $r=0.01$}
\addplot[black!60!green,mark=triangle,mark size=1.6pt,dashed] table[x=x,y=N]{generated/fig_x_2b.dat};\addlegendentry{$2k+1$, $r=1$}
\end{axis}
\end{tikzpicture}
\caption{Stability in $x$: $N(x,r)$ against $x$ for $r=0.01$ and $r=1$.}
\label{fig:x}
\end{figure}

\subsection{The exact-shape statement}\label{sec:num-shape}

\Cref{thm:main}(ii) replaces $C_fx^{1/2}Y^{3/2}$ by $\Sigma_{\MT}=\sum_{k\le y}f'(k)\MT(x;f(k))$. We evaluated $\Sigma_{\MT}$ for the rows listed in \cref{tab:shape}, together with $\sum_{k\le y}F(k)h(f(k))$, the main term of \cref{prop:single}(ii) summed over the moduli. We report
\[
N_{\MT}=\frac{\Sigma_{\MT}}{x^{1/2}Y^{3/2}},\qquad N_h=\frac{1}{Y^{3/2}}\sum_{k\le y}F(k)h(f(k)).
\]
$\MT(x;q)$ is computed from the factorisation $t=t_0t_1t_2$ of the proof of \cref{prop:MT}. For $t_0\le\sqrt{Z_1}$ the terms are summed directly. For $t_0>\sqrt{Z_1}$ we use $B(u)=u-u^2$ ($u<1$) and the closed forms
\[
A_0\sum_{(t,q)=1}\eta(t)=\frac6{\pi^2}\prod_{p\mid q}\frac{p^2-2}{p^2-1},\qquad A_0\sum_{(t,q)=1}\frac{\eta(t)}{t^2}=\frac{36}{\pi^4}\prod_{p\mid q}\frac{p^2(p^2-2)}{(p^2-1)^2},
\]
which follow from~\eqref{eq:euler} as in \cref{lem:eta}. The first $64$ terms of the subtracted partial sums are evaluated in $80$-digit arithmetic; $A_0$ and $\prod_pa_p$ are evaluated to high precision through the prime zeta function. The infinite series~\eqref{eq:MT} is thus evaluated without truncation, up to floating-point rounding. As an independent check, \cref{prop:identity} was verified with this implementation of $\MT$ against brute-force values of every other term, for $x=3000$ and all $q\le40$ (relative discrepancy at most $\SelfTest$). For $2k+1$ the number of moduli is about $x/2$, so only $x\le10^6$ was treated.

\begin{table}[htbp]
\centering
\small
\begin{tabular}{llrrrrrr}
\toprule
$f$ & $x$ & $r$ & $y$ & $N_{\mathrm{data}}$ & $N_{\mathrm{MT}}$ & $V_f/\Sigma_{\mathrm{MT}}$ & $N_h$ \\
\midrule
$k^2+1$ & $10^{6}$ & $0.01$ & 99 & 0.14185 & 0.14171 & 1.00100 & 0.13966 \\
 &  & $0.1$ & 316 & 0.15191 & 0.15191 & 0.99999 & 0.13898 \\
 &  & $1$ & 999 & 0.17136 & 0.17136 & 1.00003 & 0.13825 \\
 & $1.6\cdot 10^{7}$ & $0.01$ & 399 & 0.14117 & 0.14118 & 0.99996 & 0.13853 \\
 &  & $0.1$ & 1264 & 0.15115 & 0.15115 & 1.00003 & 0.13825 \\
 &  & $1$ & 3999 & 0.17115 & 0.17115 & 1.00000 & 0.13812 \\
 & $2.56\cdot 10^{8}$ & $0.005$ & 1131 & 0.14039 & 0.14039 & 0.99998 & 0.13829 \\
 &  & $0.01$ & 1599 & 0.14086 & 0.14086 & 0.99999 & 0.13820 \\
 &  & $0.1$ & 5059 & 0.15098 & 0.15097 & 1.00001 & 0.13812 \\
 &  & $1$ & 15999 & 0.17110 & 0.17110 & 1.00000 & 0.13809 \\
\midrule
$k^3+2$ & $10^{6}$ & $0.01$ & 21 & 0.15356 & 0.15361 & 0.99968 & 0.15544 \\
 &  & $0.1$ & 46 & 0.15503 & 0.15499 & 1.00023 & 0.14231 \\
 &  & $1$ & 99 & 0.17071 & 0.17071 & 1.00003 & 0.13890 \\
 & $1.6\cdot 10^{7}$ & $0.01$ & 54 & 0.14240 & 0.14236 & 1.00024 & 0.14132 \\
 &  & $0.1$ & 116 & 0.14975 & 0.14974 & 1.00009 & 0.13844 \\
 &  & $1$ & 251 & 0.16787 & 0.16787 & 1.00000 & 0.13708 \\
 & $2.56\cdot 10^{8}$ & $0.005$ & 108 & 0.14076 & 0.14079 & 0.99982 & 0.13851 \\
 &  & $0.01$ & 136 & 0.13999 & 0.14000 & 0.99995 & 0.13776 \\
 &  & $0.1$ & 294 & 0.14848 & 0.14848 & 1.00000 & 0.13673 \\
 &  & $1$ & 634 & 0.16671 & 0.16671 & 1.00000 & 0.13602 \\
\midrule
$2k+1$ & $2.5\cdot 10^{5}$ & $0.01$ & 1249 & 0.14376 & 0.14466 & 0.99381 & 0.14268 \\
 &  & $0.1$ & 12499 & 0.15670 & 0.15671 & 0.99992 & 0.14262 \\
 &  & $1$ & 124999 & 0.18370 & 0.18369 & 1.00002 & 0.14261 \\
 & $10^{6}$ & $0.01$ & 4999 & 0.14439 & 0.14461 & 0.99848 & 0.14263 \\
 &  & $0.1$ & 49999 & 0.15670 & 0.15670 & 0.99997 & 0.14262 \\
 &  & $1$ & 499999 & 0.18369 & 0.18369 & 0.99999 & 0.14261 \\
\bottomrule
\end{tabular}

\caption{Exact-shape comparison: the data $N_{\mathrm{data}}=N(x,r)$, the normalised sum $N_{\MT}$ of the exact main terms, the ratio $V_f/\Sigma_{\MT}$, and the normalised $h$-average $N_h$. Rows with $r\in\{0.005,0.01,0.1,1\}$ are shown; \texttt{reproduction.json} contains all $\ShapeRows$ computed rows.}
\label{tab:shape}
\end{table}

Over all $\ShapeRows$ computed rows, $V_f/\Sigma_{\MT}$ lies in $[\ShapeMin,\ShapeMax]$. At $x=2.56\cdot10^8$ it differs from $1$ by at most $\ShapeDevLargest$. So the exact main term of \cref{thm:main}(ii) accounts for the whole ratio dependence seen in \cref{tab:largest}. This agreement is much better than the proved error term $x^{5/3+\eps}+x^{1/2}Y$. Relative to $x^{1/2}Y^{3/2}$, and ignoring the unknown constant and the factor $x^\eps$, that error term is of size about $x^{-1/3}r^{-3/2}$. At $x=2.56\cdot10^8$ this is about $\XLargeInvCubeSmall$ for $r=0.01$ and $\XLargeInvCubeOne$ for $r=1$. The pair-count errors $D_x(hq)$ evidently cancel substantially when summed over $h$, which the proof does not exploit. The quantity $N_h$ approaches $C_f$ as $y$ grows, as \cref{prop:poly}(a) predicts. For $k^3+2$ the convergence is slower and less regular, because $y$ is only a few hundred and the error is of order $y^{-1/2}$. The difference $N_{\MT}-N_h$ is the secondary term of the sawtooth sums. It is small for small $r$ and grows with $r$. This is consistent with the relative bound $(Y/x)^{1/6}$ (up to $x^{\eps}$) that \cref{prop:MT} implies, but in the data the difference behaves roughly like a multiple of $\sqrt r$ (\cref{fig:shape}).

\begin{figure}[htbp]
\centering
\begin{tikzpicture}
\begin{groupplot}[group style={group size=2 by 1,horizontal sep=1.6cm},width=0.47\textwidth,height=0.4\textwidth,
  xlabel={$\sqrt r$},xmin=-0.02,xmax=1.02,ymin=0.13,ymax=0.175,tick label style={font=\footnotesize},
  yticklabel style={/pgf/number format/fixed,/pgf/number format/precision=3},legend style={font=\footnotesize},legend cell align=left]
\nextgroupplot[title={$k^2+1$, $x=2.56\cdot10^8$},ylabel={normalised value},legend pos=north west]
\addplot[only marks,mark=*,mark size=1.6pt,blue] table[x=sqrtr,y=Ndata]{generated/fig_shape_0.dat};\addlegendentry{$N_{\mathrm{data}}$}
\addplot[blue,thin] table[x=sqrtr,y=NMT]{generated/fig_shape_0.dat};\addlegendentry{$N_{\MT}$}
\addplot[black,dashed,mark=x] table[x=sqrtr,y=Nh]{generated/fig_shape_0.dat};\addlegendentry{$N_h$}
\addplot[only marks,mark=o,mark size=3pt,thick,black] table[x=sqrtr,y=N]{generated/fig_cf_0.dat};\addlegendentry{$C_f$}
\nextgroupplot[title={$k^3+2$, $x=2.56\cdot10^8$}]
\addplot[only marks,mark=square*,mark size=1.6pt,red] table[x=sqrtr,y=Ndata]{generated/fig_shape_1.dat};
\addplot[red,thin] table[x=sqrtr,y=NMT]{generated/fig_shape_1.dat};
\addplot[black,dashed,mark=x] table[x=sqrtr,y=Nh]{generated/fig_shape_1.dat};
\addplot[only marks,mark=o,mark size=3pt,thick,black] table[x=sqrtr,y=N]{generated/fig_cf_1.dat};
\end{groupplot}
\end{tikzpicture}
\caption{Data, exact main term $N_{\MT}$, averaged asymptotic main term $N_h$ and $C_f$ at the largest $x$. Both panels use the legend of the left panel: filled marks are the data, the solid line is $N_{\MT}$, the dashed line with crosses is $N_h$, and the open circle at $r=0$ is $C_f$.}
\label{fig:shape}
\end{figure}

\FloatBarrier
\subsection{Limitations}\label{sec:num-limits}

\begin{itemize}
\item The computations stop at $x=2.56\cdot10^8$. With the actual ratio $\varrho=f(y)/x$, a row lies in the range $x^{7/9+\delta}\le f(y)\le x^{1-\delta}$ of \cref{thm:main}(iii) exactly when
\[
\delta\le\delta_{\mathrm{row}}:=\min\Bigl(\frac{\log\varrho}{\log x}+\frac29,\ -\frac{\log\varrho}{\log x}\Bigr).
\]
At $x=2.56\cdot10^8$ we have $x^{-2/9}\approx\XTwoNinths$. The rows with $r\le0.01$ have $\varrho\le r<x^{-2/9}$, so $\delta_{\mathrm{row}}<0$ and they lie outside the range for every $\delta>0$. The rows with $r=1$ have $\varrho$ close to $1$ and $\delta_{\mathrm{row}}\le\DeltaRowOneMax$, so they lie in the range only for very small $\delta$; the implied constant in~(iii) depends on $\delta$ and is not explicit. Over the whole grid at this $x$, $\delta_{\mathrm{row}}$ is largest, about $\DeltaRowMax$, at $r=\DeltaRowMaxR$. For each fixed $\delta>0$, the rows with a fixed target ratio $r$ leave the range as $x$ grows, since they need $x\le\varrho^{-1/\delta}$, where $\varrho\approx r$. The values of $\delta_{\mathrm{row}}$ for all rows at the largest $x$ are recorded in \texttt{reproduction.json}.
\item The implied constants in \cref{thm:main} are not explicit, so no finite computation can confirm or refute it. The data can only be checked for consistency.
\item The extrapolated constants depend on an empirical model, and their uncertainties are heuristic envelopes, not confidence intervals. The agreement in \cref{tab:fits} is therefore qualitative evidence for the value of $C_f$, and the theorem does not rest on it.
\item The values of $C_f$ in \cref{tab:constants} come with a rigorous bound for the omitted primes, but not with a certified bound for floating-point rounding; the $40$-digit recomputation for $N=10^5$ is only a check.
\item The exact-shape comparison uses double-precision arithmetic, with $80$-digit evaluation only where cancellation occurs. Its agreement with the identity self-test is the main check of the implementation. For $2k+1$ it covers only $x\le10^6$.
\end{itemize}
